\documentclass[11pt,a4paper]{article}

\usepackage[T1]{fontenc}
\usepackage[utf8]{inputenc}
\usepackage{lmodern}
\usepackage[margin=25mm]{geometry}
\usepackage{microtype}
\usepackage{amsmath}
\usepackage{booktabs}
\usepackage{tabularx}
\usepackage{enumitem}
\usepackage{titlesec}
\usepackage[hidelinks]{hyperref}

\setlist{nosep,leftmargin=*,itemsep=2pt}
\setlength{\parindent}{0pt}
\setlength{\parskip}{6pt}
\renewcommand{\arraystretch}{1.25}
\setcounter{secnumdepth}{3}

\titleformat{\section}{\Large\bfseries}{\thesection}{0.8em}{}
\titleformat{\subsection}{\large\bfseries}{\thesubsection}{0.8em}{}
\titlespacing*{\section}{0pt}{18pt}{6pt}
\titlespacing*{\subsection}{0pt}{12pt}{4pt}

% Column types for tabularx tables.
\newcolumntype{L}{>{\raggedright\arraybackslash}X}

% Spec conventions.
\newcommand{\tbd}{\textbf{TBD}}
\newcommand{\sref}[1]{\S\ref{#1}}
\newcommand{\hex}[1]{\texttt{#1}}
\newcommand{\I}{I\textsuperscript{2}C}

\title{TH-1 Bus Replay Protocol Specification}
\author{}
\date{}

\begin{document}

\maketitle
\vspace{-2.5em}

The TH-1 bus replay protocol connects a bus capture and replay device to a host
computer. The device attaches to a bus under test, such as \I{}, samples its lines and
streams every level change to the host, where the traffic is decoded, inspected and
replayed. This document specifies how captures are encoded and streamed, how the host
configures the device, and how transports carry both.

\bigskip

\noindent\textbf{Status:} draft

\medskip

\begin{tabularx}{\linewidth}{@{}l l L@{}}
\toprule
\textbf{Version} & \textbf{Date} & \textbf{Changes} \\
\midrule
1.0 & 2026-10-19 & First version. \\
\bottomrule
\end{tabularx}

\newpage
\tableofcontents

%--------------------------------------------------------------------
\section{Scope}\label{sec:scope}

In scope:
\begin{itemize}
  \item The channel model, and how logical channels map onto data channels.
  \item The time model, records and heartbeats on the device $\to$ host data path.
  \item Loss detection and reporting.
  \item Control messages: hello, registers, configuration, loss report, capture
        lifecycle.
  \item Transports: the contract every transport must meet, and the bindings this
        document defines.
\end{itemize}

Out of scope:
\begin{itemize}
  \item The device's internal design (sampling method, buffer implementation).
  \item The host interface and its data model.
  \item Protocol decoding (\I{} or otherwise). Decode is a host-side step over the raw
        pin states this protocol delivers.
\end{itemize}

%--------------------------------------------------------------------
\section{Terminology}\label{sec:terms}

\begin{tabularx}{\linewidth}{@{}l L@{}}
\toprule
\textbf{Term} & \textbf{Meaning} \\
\midrule
Device & The capture/replay hardware attached to the bus under test. \\
Host & The PC running the analyser interface. \\
Logical channel & One device pin, with a channel number fixed by the device, enabled in the configuration. The unit the protocol reasons about. \\
Data channel & A transport-provided byte stream that carries pin states. \\
Control channel & The single transport-provided message channel for everything that is not pin state. \\
Tick & One sample period. Global, shared by all logical channels. \\
Record & A level change on one logical channel (\sref{sec:record}). \\
Heartbeat & A periodic entry, not a record, that marks the same instant on every logical channel (\sref{sec:heartbeat}). \\
Gap & The number of ticks between a record and its reference point (\sref{sec:gap}). \\
Register & A device setting with an ID, written by the host (\sref{sec:registers}). \\
Grid point & A tick at which a heartbeat is due: every $C$ ticks from the start of the capture (\sref{sec:emission}). \\
Position & A tick on the capture timeline, as grid count and offset (\sref{sec:position}). \\
Entry & A record or a heartbeat on a data channel. \\
Binding & The part of a transport that carries one data channel, or the control channel, over one medium (\sref{sec:transports}). \\
Untrusted span & A span of a data channel the host cannot trust after transport loss (\sref{sec:transportloss}). \\
\bottomrule
\end{tabularx}

\medskip

\begin{tabularx}{\linewidth}{@{}l L@{}}
\toprule
\textbf{Symbol} & \textbf{Meaning} \\
\midrule
$N$ & The number of channels the device has (\sref{sec:hello}). \\
$n$ & The number of enabled logical channels (\sref{sec:registers}). \\
$D$ & The number of data channels (\sref{sec:channels}). \\
$C$ & The heartbeat interval in ticks (\sref{sec:interval}). \\
$W$ & The width of a record in bytes, per data channel (\sref{sec:record}). \\
$t$ & The width of a record's tag in bits, per data channel (\sref{sec:record}). \\
\bottomrule
\end{tabularx}

%--------------------------------------------------------------------
\section{Design principles}\label{sec:principles}

\begin{enumerate}
  \item \textbf{Streamed.} Captures flow to the host while the capture runs, not as a
        buffer download after it stops.
  \item \textbf{Lightweight.} Minimal bytes on the wire and minimal device-side work.
        The host may spend processing freely.
  \item \textbf{Separate control channel.} Everything that is not pin state goes on it.
  \item \textbf{Data channels carry only pin states}, in either direction.
  \item \textbf{Transport agnostic.} The core works over any transport that meets the
        contract (\sref{sec:contract}).
  \item \textbf{Bus-protocol agnostic.} Nothing in the protocol may assume the channel
        count, clock rate or traffic density of a particular bus protocol.
  \item \textbf{Fixed sample rate}, set in the configuration and constant for a
        capture.
  \item \textbf{Encoded.} Captures are per-channel change records, not raw fixed-rate
        samples.
  \item \textbf{Buffered.} The device has one buffer and one tick counter. Capture may
        outrun the transport; the backlog drains while the bus is idle.
\end{enumerate}

%--------------------------------------------------------------------
\section{Channel model}\label{sec:channels}

The protocol defines \textbf{one control channel and $D$ data channels}. Every
transport must provide both. How it does so, and the value of $D$, belong to the
transport (\sref{sec:transports}).

\subsection{Control channel}\label{sec:control}

\begin{itemize}
  \item Two-way, low rate.
  \item Message-framed and integrity-checked by its control binding
        (\sref{sec:controlbindings}): a message arrives whole and intact, or not at
        all.
  \item Delivery is not guaranteed.
  \item Device $\to$ host: hello (\sref{sec:hello}), configuration replies
        (\sref{sec:config}), loss reports (\sref{sec:lossreport}).
  \item Host $\to$ device: configuration (\sref{sec:config}).
  \item Messages: \sref{sec:messages}.
\end{itemize}

\subsection{Data channels}\label{sec:data}

\begin{itemize}
  \item Each data channel is an ordered byte stream in each direction.
  \item They carry pin states and nothing else.
    \begin{itemize}
      \item Device $\to$ host: records (\sref{sec:time}) and heartbeats
            (\sref{sec:heartbeat}).
      \item Host $\to$ device: replay (\sref{sec:replay}).
    \end{itemize}
  \item The core protocol adds no escapes or framing to a data channel: entries follow
        one another directly. Any framing a transport needs belongs to that transport.
  \item $D$ is set by the transport and must not limit how many logical channels exist.
  \item The protocol sets no maximum number of logical channels. The limit is the
        device's.
  \item Which logical channels each data channel carries is set by the mapping
        (\sref{sec:mapping}).
\end{itemize}

\subsection{Mapping}\label{sec:mapping}

This defines the device $\to$ host direction. The host $\to$ device direction is
\tbd{} (\sref{sec:replay}).

\begin{itemize}
  \item The transport numbers its data channels $0 \ldots D-1$. The hello carries $D$
        (\sref{sec:hello}).
  \item The logical channels, in ascending channel number, are split into $D$
        contiguous blocks, one per data channel, in data channel order. A record
        carries its index within the block (\sref{sec:record}).
  \item The split register (\sref{sec:registers}) sets where the blocks split. Empty,
        its default, means automatic: the blocks are $\lfloor n/D \rfloor$ channels,
        and the first $n \bmod D$ blocks have one more.
  \item Otherwise, a channel in the split set starts a new block. It need not be
        enabled: the split is by channel number.
  \item A block may be empty, for example with fewer than $D-1$ channels in the split
        set, or none enabled between two of them. Its data channel then carries no
        records, only heartbeats (\sref{sec:emission}).
\end{itemize}

For example, with channels 0, 1, 2, 3, 4, 6 and 7 enabled and $D = 3$:

\begin{tabularx}{\linewidth}{@{}l l l L@{}}
\toprule
\textbf{Split} & \textbf{Data channel 0} & \textbf{Data channel 1} & \textbf{Data channel 2} \\
\midrule
empty & 0, 1, 2 & 3, 4 & 6, 7 \\
\{2, 6\} & 0, 1 & 2, 3, 4 & 6, 7 \\
\{6\} & 0, 1, 2, 3, 4 & 6, 7 & none \\
\bottomrule
\end{tabularx}

\subsection{Contract}\label{sec:contract}

Across its data channels, a transport must guarantee:
\begin{itemize}
  \item On each data channel, records and heartbeats arrive in the order the device
        wrote them.
  \item Every heartbeat arrives intact on every device $\to$ host data channel
        (\sref{sec:emission}), or the transport detects that it did not
        (\sref{sec:transportloss}).
\end{itemize}

Records carry no integrity guarantee.

Ordering \emph{between} logical channels carried on different data channels is not
guaranteed. The host aligns them using heartbeats (\sref{sec:hostuse}).

A transport must also give the host a way to reset the device (\sref{sec:lifecycle}).

%--------------------------------------------------------------------
\section{Time and records}\label{sec:time}

\subsection{Tick}\label{sec:tick}

One tick is one sample period. The sample rate is fixed by the configuration for the
whole capture. There is one tick counter on the device, shared by all logical channels.

The device samples every tick. It does not start a capture at a sample rate it cannot
sustain (\sref{sec:config}).

\subsection{Record}\label{sec:record}

A record is emitted for every sampled level change on a logical channel. Nothing is
filtered. Logically it is:

\begin{center}
\texttt{\{ channel, new level, gap \}}
\end{center}

On the wire it is one value of $W$ bytes, low byte first:

\begin{itemize}
  \item $v = +\mathrm{gap}$ for new level 1, $v = -\mathrm{gap}$ for new level 0.
  \item $z = (v \ll 1) \oplus (v \gg 31)$, with an arithmetic shift (ZigZag). Bit 0 of
        $z$ is 1 for level 0, and 0 for level 1.
  \item \textbf{Tag.} The record's channel is identified by its index within the
        data channel's block (\sref{sec:mapping}), counted from 0 in ascending channel
        number, in $t = \lceil \log_2 b \rceil$ bits, where $b$ is the block size. A
        data channel carrying one logical channel or none has $t = 0$: no tag.
  \item The value is $(z \ll t) \mid \mathrm{index}$.
  \item $W$ is the smallest number of bytes that holds $(2C + 1) \times 2^t - 1$. Every
        record on a data channel is $W$ bytes.
\end{itemize}

$W$ and $t$ are per data channel. Both sides derive them from the configuration.

\subsection{Gap}\label{sec:gap}

A record's gap is the number of ticks since the more recent of:
\begin{itemize}
  \item the previous record on the same logical channel, or
  \item the most recent heartbeat written before it (\sref{sec:heartbeat}).
\end{itemize}

A heartbeat resets the gap reference of every logical channel at once. A gap is at
least 1.

A record that was dropped (\sref{sec:devicebehaviour}) was never written and resets
nothing.

A gap never exceeds $C$ (\sref{sec:pending}).

%--------------------------------------------------------------------
\section{Heartbeat}\label{sec:heartbeat}

A periodic marker of the same instant on every logical channel. It re-anchors the
timeline so that alignment between channels and recovery from errors do not depend on
summed gaps alone.

\subsection{Content}\label{sec:hbcontent}

\begin{itemize}
  \item A heartbeat is one entry for all logical channels, not a record. It is marked
        by a value of 0. No record has that value, since a record's $z$ is at least 1
        (\sref{sec:record}, \sref{sec:gap}).
  \item It carries the level of every logical channel as a bitmask, one bit per
        logical channel. A heartbeat is therefore a keyframe for all channels in one
        write.
  \item The bitmask holds each channel's level at the end of the heartbeat's tick,
        including any change on that tick.
  \item It carries a grid index: the number of its grid point, counted from the start
        of the capture, mod 256.
\end{itemize}

On the wire:

\begin{tabularx}{\linewidth}{@{}l l L@{}}
\toprule
\textbf{Field} & \textbf{Size} & \textbf{Contents} \\
\midrule
Marker & $W$ bytes (\sref{sec:record}) & all zero, with no tag \\
Grid index & 1 byte & grid index \\
Bitmask & $\lceil n/8 \rceil$ bytes & the $i$-th logical channel in ascending channel number, counted from 0, is bit $(i \bmod 8)$ of byte $\lfloor i/8 \rfloor$. Unused bits are 0. \\
\bottomrule
\end{tabularx}

A heartbeat is $W + 1 + \lceil n/8 \rceil$ bytes, with the $W$ of the data channel it
is on. It is fixed for the capture.

\subsection{Emission}\label{sec:emission}

\begin{itemize}
  \item A capture begins with a heartbeat.
  \item A heartbeat marks tick 0 at its grid point, even if written later
        (\sref{sec:pending}), and resets every gap reference (\sref{sec:gap}).
  \item A heartbeat is due at every grid point.
  \item Within a tick, records are written before the heartbeat.
  \item Heartbeats are global: the device writes one for all logical channels. It goes
        on every device $\to$ host data channel, with that data channel's $W$, so each
        stream can anchor itself.
\end{itemize}

\subsection{Interval}\label{sec:interval}

$C$ is set by the configuration (\sref{sec:registers}), within the range the device
lists in its hello (\sref{sec:hello}). The protocol fixes no default or range.

\subsection{Pending heartbeat}\label{sec:pending}

A heartbeat is never skipped.

\begin{itemize}
  \item If the buffer has no room for a heartbeat at its grid point, it is pending. The
        device keeps its grid index and the levels at the grid point.
  \item A pending heartbeat is written as soon as there is room, before any record.
        Records that arrive meanwhile and do not fit after it are dropped
        (\sref{sec:devicebehaviour}). Records are never held back.
  \item The device does not start a capture with a $C$ shorter than the time needed to
        send one heartbeat on every data channel (\sref{sec:config}). While a
        heartbeat is pending nothing else is written and the buffer drains, so it is
        always written before the next grid point.
  \item A transport that stops draining is outside the protocol.
\end{itemize}

Every record therefore follows the heartbeat of the grid point before it, and a gap
never exceeds $C$.

\subsection{Position}\label{sec:position}

A position is a tick on the capture timeline, sent where a message needs to name one
(\sref{sec:lossreportmsg}).

\begin{tabularx}{\linewidth}{@{}l l L@{}}
\toprule
\textbf{Field} & \textbf{Size} & \textbf{Contents} \\
\midrule
Grid count & 4 bytes & the number of the grid point at or before the tick, counted from the start of the capture \\
Offset & 4 bytes & ticks since that grid point, $0 \ldots C-1$ \\
\bottomrule
\end{tabularx}

\begin{itemize}
  \item The tick is grid count $\times\, C\, +$ offset.
  \item The grid count is mod $2^{32}$. The host unwraps it against the grid points it
        has seen.
  \item A heartbeat's grid index is its grid count mod 256 (\sref{sec:hbcontent}).
  \item A position names a tick, not an epoch. A tick that is a grid point always has
        offset 0, even for a record on that tick whose gap counts from the previous
        heartbeat (\sref{sec:emission}).
\end{itemize}

\subsection{Host use}\label{sec:hostuse}

\begin{itemize}
  \item \textbf{Alignment.} Heartbeats with the same grid index mark the same instant
        on every data channel. The host aligns data channels on that.
  \item \textbf{Known extent.} A logical channel's timeline is known up to its latest
        record or heartbeat. A quiet channel is therefore known to within one interval
        $C$.
  \item \textbf{Epoch length.} The time between two heartbeats received on a data
        channel is their grid index difference (mod 256) $\times\, C$ ticks. This holds
        across heartbeats lost in transport, and is exact unless 256 or more intervals
        pass between two received heartbeats.
\end{itemize}

%--------------------------------------------------------------------
\section{Loss}\label{sec:loss}

Loss is never repaired. There are two kinds:
\begin{itemize}
  \item \textbf{Device loss.} The buffer fills when capture outruns the transport. The
        device detects the loss itself and reports it on the control channel
        (\sref{sec:devicebehaviour}--\sref{sec:hostbehaviour}).
  \item \textbf{Transport loss.} Bytes lost or corrupted on a data channel between
        device and host (\sref{sec:transportloss}).
\end{itemize}

\subsection{Device behaviour}\label{sec:devicebehaviour}

\begin{itemize}
  \item A record is written only if the buffer has room for it ($W$ bytes,
        \sref{sec:record}). Otherwise it is dropped whole, never partly written, and
        is counted (\sref{sec:lossreport}).
  \item While a heartbeat is pending, a record is written only after it, and only if
        both fit (\sref{sec:pending}).
  \item Nothing is sent on the data path. Gaps after a drop stay correct
        (\sref{sec:gap}).
\end{itemize}

\subsection{Loss report}\label{sec:lossreport}

A report says which channels lost records and how many (the counters), and when (the
loss start and end).

The device keeps a \textbf{lost-records} counter for each logical channel: the number
of records on that channel dropped for lack of buffer room. Each counter:
\begin{itemize}
  \item is an unsigned 32-bit integer,
  \item only increases, and wraps at $2^{32}$,
  \item is never cleared except by a device reset (\sref{sec:lifecycle}).
\end{itemize}

Drops are grouped into losses, and the device reports each loss on the control channel
(\sref{sec:lossreportmsg}):
\begin{itemize}
  \item \textbf{Start.} A loss starts when a record is dropped while no loss is open.
        The position (\sref{sec:position}) of that drop is the loss start.
  \item \textbf{End.} The position of the latest drop in the loss is the loss end. The
        loss is over when the next entry, a record or a heartbeat, is written.
        Heartbeats are never skipped (\sref{sec:pending}), so a loss is over within
        one interval of the buffer having room, even on an idle bus.
  \item \textbf{Report.} When a loss is over, a report is due. It carries the loss
        start, the loss end and every counter's value at that point. The device sends
        it when the control channel can take it.
  \item \textbf{Merging.} If the previous report has not been sent when a loss is
        over, the loss is merged into it: the report keeps its loss start and takes
        the new loss end and counters. At most one report waits at a time.
  \item \textbf{Sequence.} Each report carries a sequence number. It is 0 for the
        first report after a boot and goes up by one for each report sent.
\end{itemize}

\subsection{Host behaviour}\label{sec:hostbehaviour}

\begin{itemize}
  \item \textbf{Loss report.} Report an overrun. For each logical channel whose counter
        rose, mark possible lost transitions from the loss start to the loss end.
        Records inside that range are kept. Other channels are not marked. Do not
        repair or guess lost transitions.
  \item \textbf{Counters.} The host keeps the last value of each counter and takes the
        unsigned difference, so a wrap reads correctly. After a reset it starts again
        from 0.
  \item \textbf{Sequence skip.} A report was lost. The counter differences are still
        exact. For each channel whose counter rose, mark from the loss end of the
        previous report received to the loss end of this one, or from the start of the
        capture if none was received since the reset.
  \item \textbf{Merged range.} A merged report's range can include records that were
        written. It may mark more than was lost, never less.
  \item \textbf{Repeated level.} After a dropped record, the next record on that
        channel may carry the channel's previous level. The host accepts it.
\end{itemize}

\subsection{Transport loss}\label{sec:transportloss}

The transport must detect loss or corruption of heartbeats (\sref{sec:contract}). It
may also detect it for records.

\begin{itemize}
  \item When the transport detects loss or corruption on a data channel, the host marks
        an \textbf{untrusted span} on every logical channel carried on that data
        channel, from the last point the transport confirmed good to the point the loss
        was detected.
  \item Levels and gaps are trusted again from the next heartbeat the transport
        confirms good.
  \item Heartbeats lost inside the span are accounted for by the grid index
        (\sref{sec:hostuse}).
\end{itemize}

%--------------------------------------------------------------------
\section{Control messages}\label{sec:messages}

Every control message has the same shape. The transport delivers it whole
(\sref{sec:control}), so its length is known.

\begin{tabularx}{\linewidth}{@{}l l L@{}}
\toprule
\textbf{Field} & \textbf{Size} & \textbf{Contents} \\
\midrule
Type & 1 byte & which message \\
Body & rest of the message & depends on the type \\
\bottomrule
\end{tabularx}

\medskip

\begin{tabularx}{\linewidth}{@{}l l L@{}}
\toprule
\textbf{Type} & \textbf{Direction} & \textbf{Message} \\
\midrule
\hex{0x01} & device $\to$ host & hello (\sref{sec:hello}) \\
\hex{0x02} & host $\to$ device & configuration (\sref{sec:config}) \\
\hex{0x03} & device $\to$ host & configuration reply (\sref{sec:config}) \\
\hex{0x04} & device $\to$ host & loss report (\sref{sec:lossreportmsg}) \\
\bottomrule
\end{tabularx}

\begin{itemize}
  \item Integers are unsigned and sent low byte first.
  \item \textbf{str}: a 1-byte length, then that many ASCII bytes, with no
        terminator. It may be empty.
  \item \textbf{bytes}: a 1-byte length, then that many raw bytes.
  \item A receiver ignores a message of a type it does not know.
\end{itemize}

\subsection{Hello}\label{sec:hello}

Sent by the device on every boot (power-on or reset), device $\to$ host. It tells the
host the device is ready.

\begin{tabularx}{\linewidth}{@{}l l L@{}}
\toprule
\textbf{Field} & \textbf{Size} & \textbf{Contents} \\
\midrule
Type & 1 byte & \hex{0x01} \\
Version major & 1 byte & major part of this document's version (change history above). The host refuses one it does not know. \\
Version minor & 1 byte & minor part of this document's version. The host accepts a newer one. \\
Model & str & device model \\
Firmware & str & firmware version \\
Unique ID & bytes & an ID unique to this unit \\
Buffer size & 4 bytes & capture buffer size in bytes \\
Channel count & 2 bytes & $N$ \\
Data channel count & 1 byte & $D$ (\sref{sec:mapping}) \\
Channel names & $N \times$ str & one per channel, in channel number order \\
Register count & 1 byte & number of register entries \\
Registers & per entry & one entry per register (\sref{sec:registers}), below \\
\bottomrule
\end{tabularx}

\medskip

Register entry:

\begin{tabularx}{\linewidth}{@{}l l L@{}}
\toprule
\textbf{Field} & \textbf{Size} & \textbf{Contents} \\
\midrule
ID & 1 byte & register ID \\
Type & 1 byte & type code \\
Flags & 1 byte & bit 0: writable. Other bits are 0. \\
Name & str & short name \\
Unit & str & unit; may be empty \\
Allowed values & kind and data & the register's allowed values (\sref{sec:registers}) \\
Current value & per the type & after a boot, the default \\
\bottomrule
\end{tabularx}

There is no query message. To read the device's state, the host resets it and reads the
hello.

\subsection{Registers}\label{sec:registers}

Every setting is a register with a 1-byte ID. IDs \hex{0x00}--\hex{0x7F} are
reserved for the core, which defines the registers below with fixed IDs, the same on
every device. IDs \hex{0x80}--\hex{0xFF} are for registers a device or transport adds;
the host learns them from the hello.

\begin{tabularx}{\linewidth}{@{}l l l l l L@{}}
\toprule
\textbf{ID} & \textbf{Register} & \textbf{Type} & \textbf{Unit} & \textbf{Default} & \textbf{Allowed values} \\
\midrule
\hex{0x00} & capture & bool & none & off & list: on. Only while off (\sref{sec:lifecycle}) \\
\hex{0x01} & sample rate & u32 & Hz & device & device, as a range, list or divider \\
\hex{0x02} & $C$ & u32 & ticks & device & device, as a range, minimum at least 1 \\
\hex{0x03} & enabled channels & channel set & none & device & any \\
\hex{0x04} & split & channel set & none & empty & any \\
\bottomrule
\end{tabularx}

\medskip

\begin{tabularx}{\linewidth}{@{}l l l L@{}}
\toprule
\textbf{Code} & \textbf{Type} & \textbf{Size} & \textbf{Encoding} \\
\midrule
\hex{0x00} & bool & 1 byte & \hex{00} off, \hex{01} on \\
\hex{0x01} & u32 & 4 bytes & unsigned, low byte first \\
\hex{0x02} & channel set & $\lceil N/8 \rceil$ bytes & Channel $i$ is bit $(i \bmod 8)$ of byte $\lfloor i/8 \rfloor$; 1 means in the set. Unused bits are 0. \\
\bottomrule
\end{tabularx}

Other type codes are reserved.

Allowed values, as the hello lists them (\sref{sec:hello}), start with a kind byte.
Values are in the register's type.

\begin{tabularx}{\linewidth}{@{}l l L@{}}
\toprule
\textbf{Kind} & \textbf{Meaning} & \textbf{Followed by} \\
\midrule
\hex{0x00} & any value of the type & nothing \\
\hex{0x01} & range, inclusive & min, max \\
\hex{0x02} & list & count (u16), then the values \\
\hex{0x03} & divider & base, min divisor, max divisor (each u32) \\
\bottomrule
\end{tabularx}

\begin{itemize}
  \item \textbf{Divider.} The allowed values are base $\div\, d$ for every whole $d$ from
        min to max divisor for which the result is whole. The min divisor is at least
        1. For u32 registers only.
  \item \textbf{Sample rate.} The device accepts only rates it can produce exactly; it
        does not round.
  \item \textbf{Channel sets.} A bit for a channel the device does not have is
        rejected.
  \item \textbf{Enabled channels.} At least one channel must be enabled to start a
        capture (\sref{sec:config}).
\end{itemize}

\subsection{Configuration}\label{sec:config}

The host writes one register per configuration message, and sends the next only after
the reply to the last.

\begin{tabularx}{\linewidth}{@{}l l L@{}}
\toprule
\textbf{Field} & \textbf{Size} & \textbf{Contents} \\
\midrule
Type & 1 byte & \hex{0x02} \\
Register ID & 1 byte & the register (\sref{sec:registers}) \\
Value & per the register's type & the new value \\
\bottomrule
\end{tabularx}

The device replies accepted, or rejected with a reason:

\begin{tabularx}{\linewidth}{@{}l l L@{}}
\toprule
\textbf{Field} & \textbf{Size} & \textbf{Contents} \\
\midrule
Type & 1 byte & \hex{0x03} \\
Register ID & 1 byte & the register written \\
Status & 1 byte & \hex{0x00} accepted, otherwise a reason code \\
Message & str & text from the device; may be empty \\
Value & per the type, rest of the message & the register's value after the write: the new value if accepted, the unchanged value if rejected, empty if the register is unknown \\
\bottomrule
\end{tabularx}

\medskip

\begin{tabularx}{\linewidth}{@{}l L@{}}
\toprule
\textbf{Code} & \textbf{Reason} \\
\midrule
\hex{0x01} & unknown register \\
\hex{0x02} & register not writable \\
\hex{0x03} & wrong value length for the type, or message too short \\
\hex{0x04} & value not in the allowed values \\
\hex{0x05} & capture already started \\
\hex{0x06} & start: sample rate not sustainable \\
\hex{0x07} & start: $C$ too short to send one heartbeat on every data channel \\
\hex{0x08} & start: no channel enabled \\
\hex{0x09} & start: split set holds more than $D-1$ channels \\
\hex{0x80}--\hex{0xFF} & device-specific; the message explains it \\
\bottomrule
\end{tabularx}

A start rejection (\hex{0x06}--\hex{0x09}) is a reply to the capture register,
\hex{0x00}.

\begin{itemize}
  \item Each write is checked on its own: the register exists and is writable, the
        value has the right type and is within its allowed values, and capture has not
        started. If a check fails, the write is rejected and nothing changes.
  \item If no reply arrives (\sref{sec:control}), the host may send the same write
        again. Every write can be repeated safely: a repeated start after one that
        succeeded is answered \hex{0x05}, which means capture is on.
  \item Combined limits are checked when capture is started: the sample rate must be
        one the device can sustain (\sref{sec:tick}), $C$ must be long enough to send
        one heartbeat on every data channel (\sref{sec:pending}), at least one channel
        must be enabled, and the split set may hold at most $D-1$ channels
        (\sref{sec:mapping}). If they fail, the start is rejected and capture stays
        off. Before start, registers may be written in any order.
\end{itemize}

\subsection{Loss report}\label{sec:lossreportmsg}

Sent by the device after a loss (\sref{sec:lossreport}).

\begin{tabularx}{\linewidth}{@{}l l L@{}}
\toprule
\textbf{Field} & \textbf{Size} & \textbf{Contents} \\
\midrule
Type & 1 byte & \hex{0x04} \\
Sequence & 4 bytes & report sequence number \\
Loss start & 8 bytes & position (\sref{sec:position}) of the first drop \\
Loss end & 8 bytes & position of the latest drop \\
Lost records & $n \times 4$ bytes & one counter per logical channel, in ascending channel number \\
\bottomrule
\end{tabularx}

\subsection{Capture lifecycle}\label{sec:lifecycle}

\begin{itemize}
  \item There are no triggers.
  \item Capture starts when the host sets the capture register (\sref{sec:registers}) to
        on. This is one-way: it cannot be set back to off.
  \item On an accepted start, the device sends the reply before it writes the first
        heartbeat.
  \item Once capture has started, the device rejects every configuration message.
  \item Capture stops when the host resets the device. This is outside the protocol:
        buffered data is discarded, and the device returns to its power-on state with
        capture off and must be configured again.
\end{itemize}

\subsection{Replay}\label{sec:replay}

Replay travels host $\to$ device on the data channels. Its format is \tbd{}.

%--------------------------------------------------------------------
\section{Transports}\label{sec:transports}

A transport is made of \textbf{bindings}. A data binding carries one data channel over
one medium. A control binding carries the control channel over one medium. A device
combines one control binding with $D$ data channels, each on a data binding, and a way
to reset the device (\sref{sec:contract}).

\begin{itemize}
  \item Bindings combine freely. The combination is not announced: the host knows it
        before the hello.
  \item A binding setting the host must change is a device register
        (\hex{0x80}--\hex{0xFF}, \sref{sec:registers}).
  \item Everything else (pins, rates, modes, endpoints) is implementation.
\end{itemize}

\subsection{Data bindings}\label{sec:databindings}

A data binding defines:
\begin{itemize}
  \item how the data channel's bytes are carried, in order (\sref{sec:contract});
  \item how loss or corruption of a heartbeat is detected (\sref{sec:contract},
        \sref{sec:transportloss}).
\end{itemize}

\subsubsection{Serial (UART)}\label{sec:serial}

\begin{itemize}
  \item One UART is one data channel. The device's TX carries device $\to$ host; its
        RX carries host $\to$ device (\sref{sec:replay}).
  \item A \textbf{check byte} follows every heartbeat on device $\to$ host. It is a
        CRC-8 over every byte sent since the previous check byte, up to and including
        this heartbeat; for the first heartbeat, from the start of the capture.
  \item CRC-8: polynomial \hex{0x07}, initial value \hex{0xFF}, MSB first, no
        reflection, no final XOR. Reset after each check byte.
  \item A mismatch marks an untrusted span (\sref{sec:transportloss}) from the
        previous good check byte to this one.
  \item The stream has no other resynchronisation aid. After a mismatch the host finds
        the next heartbeat itself; how is host implementation.
  \item A heartbeat on serial is $W + 1 + \lceil n/8 \rceil + 1$ bytes, counting its
        check byte. The start check on $C$ (\sref{sec:pending}) counts it too.
\end{itemize}

\subsection{Control bindings}\label{sec:controlbindings}

A control binding defines:
\begin{itemize}
  \item how a message arrives whole, with its length known (\sref{sec:control},
        \sref{sec:messages});
  \item how integrity is checked (\sref{sec:control});
  \item how the device sends a message unprompted (hello, loss report);
  \item the maximum message length.
\end{itemize}

\subsubsection{SPI}\label{sec:spi}

\begin{itemize}
  \item SCLK, MOSI, MISO and CS are normal SPI: the host is master, the device is
        slave.
  \item \textbf{IRQ} is driven by the device, active low, and pulled up by the host,
        so an unpowered or reset device reads as nothing waiting. The device asserts
        it while it has a message to send, and the host then runs a transaction.
  \item The host may start a transaction at any time to send its own message. It does
        not poll.
\end{itemize}

A transaction runs from CS low to CS high. Both lines carry the same layout at once:
MOSI the host's message, MISO the device's.

\begin{tabularx}{\linewidth}{@{}l l L@{}}
\toprule
\textbf{Field} & \textbf{Size} & \textbf{Contents} \\
\midrule
Length & 2 bytes & length $L$ of this side's message, low byte first; 0 for no message \\
Message & $L$ bytes & the control message (\sref{sec:messages}) \\
Check & 2 bytes & CRC-16 over Length and Message; absent when $L = 0$ \\
\bottomrule
\end{tabularx}

\begin{itemize}
  \item The host reads the device's Length first and clocks until both sides'
        messages and checks are through, padding its own side with \hex{00}.
  \item A device Length of \hex{0xFFFF} means the device was not ready. Nothing moved
        in either direction, and the host retries.
  \item One message per side per transaction. IRQ stays asserted while more are
        queued.
  \item The maximum message length is \hex{0xFFFE} bytes.
  \item CRC-16: polynomial \hex{0x1021}, initial value \hex{0xFFFF}
        (CRC-16/CCITT-FALSE). A receiver discards a message whose check fails.
  \item A message is sent once its transaction ends. There are no acknowledgements
        (\sref{sec:control}).
\end{itemize}

%--------------------------------------------------------------------
\section{Open items}\label{sec:open}

\subsection*{Replay}

\begin{tabularx}{\linewidth}{@{}r L l@{}}
\toprule
\textbf{\#} & \textbf{Item} & \textbf{Section} \\
\midrule
8 & Replay & \sref{sec:replay} \\
\bottomrule
\end{tabularx}

\subsection*{Transports}

\begin{tabularx}{\linewidth}{@{}r L l@{}}
\toprule
\textbf{\#} & \textbf{Item} & \textbf{Section} \\
\midrule
11 & USB data and control bindings & \sref{sec:transports} \\
12 & Serial control binding & \sref{sec:transports} \\
\bottomrule
\end{tabularx}

\end{document}
